\documentclass{amsart}
% For dealing with references we use the comment environment
\usepackage{verbatim}
\newenvironment{reference}{\comment}{\endcomment}
%\newenvironment{reference}{}{}
\newenvironment{slogan}{\comment}{\endcomment}
\newenvironment{history}{\comment}{\endcomment}

% For commutative diagrams we use Xy-pic
\usepackage[all]{xy}

% We use 2cell for 2-commutative diagrams.
\xyoption{2cell}
\UseAllTwocells

% We use multicol for the list of chapters between chapters
\usepackage{multicol}

% This is generally recommended for better output
\usepackage{lmodern}
\usepackage[T1]{fontenc}

% For cross-file-references

% Package for hypertext links:
\usepackage{hyperref}

% For any local file, say "hello.tex" you want to link to please

% Theorem environments.
%
\theoremstyle{plain}
\newtheorem{theorem}[subsection]{Theorem}
\newtheorem{proposition}[subsection]{Proposition}
\newtheorem{lemma}[subsection]{Lemma}

\theoremstyle{definition}
\newtheorem{definition}[subsection]{Definition}
\newtheorem{example}[subsection]{Example}
\newtheorem{exercise}[subsection]{Exercise}
\newtheorem{situation}[subsection]{Situation}

\theoremstyle{remark}
\newtheorem{remark}[subsection]{Remark}
\newtheorem{remarks}[subsection]{Remarks}

\numberwithin{equation}{subsection}

% Macros
%
\def\lim{\mathop{\mathrm{lim}}\nolimits}
\def\colim{\mathop{\mathrm{colim}}\nolimits}
\def\Spec{\mathop{\mathrm{Spec}}}
\def\Hom{\mathop{\mathrm{Hom}}\nolimits}
\def\Ext{\mathop{\mathrm{Ext}}\nolimits}
\def\SheafHom{\mathop{\mathcal{H}\!\mathit{om}}\nolimits}
\def\SheafExt{\mathop{\mathcal{E}\!\mathit{xt}}\nolimits}
\def\Sch{\mathit{Sch}}
\def\Mor{\mathop{\mathrm{Mor}}\nolimits}
\def\Ob{\mathop{\mathrm{Ob}}\nolimits}
\def\Sh{\mathop{\mathit{Sh}}\nolimits}
\def\NL{\mathop{N\!L}\nolimits}
\def\CH{\mathop{\mathrm{CH}}\nolimits}
\def\proetale{{pro\text{-}\acute{e}tale}}
\def\etale{{\acute{e}tale}}
\def\QCoh{\mathit{QCoh}}
\def\Ker{\mathop{\mathrm{Ker}}}
\def\Im{\mathop{\mathrm{Im}}}
\def\Coker{\mathop{\mathrm{Coker}}}
\def\Coim{\mathop{\mathrm{Coim}}}

% Boxtimes
%
\DeclareMathSymbol{\boxtimes}{\mathbin}{AMSa}{"02}

%
% Macros for moduli stacks/spaces
%
\def\QCohstack{\mathcal{QC}\!\mathit{oh}}
\def\Cohstack{\mathcal{C}\!\mathit{oh}}
\def\Spacesstack{\mathcal{S}\!\mathit{paces}}
\def\Quotfunctor{\mathrm{Quot}}
\def\Hilbfunctor{\mathrm{Hilb}}
\def\Curvesstack{\mathcal{C}\!\mathit{urves}}
\def\Polarizedstack{\mathcal{P}\!\mathit{olarized}}
\def\Complexesstack{\mathcal{C}\!\mathit{omplexes}}
% \Pic is the operator that assigns to X its picard group, usage \Pic(X)
% \Picardstack_{X/B} denotes the Picard stack of X over B
% \Picardfunctor_{X/B} denotes the Picard functor of X over B
\def\Pic{\mathop{\mathrm{Pic}}\nolimits}
\def\Picardstack{\mathcal{P}\!\mathit{ic}}
\def\Picardfunctor{\mathrm{Pic}}
\def\Deformationcategory{\mathcal{D}\!\mathit{ef}}

\setlength{\emergencystretch}{3em}
\begin{document}
\title[Stacks Project, Tag 0E3F]{574 --- Subcurve degree bounds imply global generation of line bundles on proper reduced connected curves}
\maketitle
\noindent Copyright (C) 2005--2025 Johan de Jong. Stacks Project authors. Source: \href{https://stacks.math.columbia.edu/tag/0E3F}{Tag 0E3F}.

\noindent Editorial difficulty note (not author text). Difficulty H2: A failure of global generation must be converted by duality into a common component on which every dual map vanishes. The proof uses an infinite-field vector-space argument and component induction to factor all maps through the complementary subcurve, then derives incompatible injectivity and surjectivity properties of the boundary cohomology maps..

\noindent Editorial provenance note (not author text). History: excerpt from revision \texttt{a0444\allowbreak 6e57e\allowbreak c1fbc\allowbreak 252a8\allowbreak 71afc\allowbreak ec775\allowbreak 2fb28\allowbreak 07b14}, curves.tex, lines 5537--5702. Reference presentation was adapted; original-author context and core support blocks were retained or added. The mathematical wording is unchanged. Licensed under GNU FDL 1.2 or later; no invariant sections or cover texts. See ../licenses/COPYING. Context blocks reproduce author definitions and supporting results; they are not additional counted problems.

\begin{lemma}
\label{lemma-duality-dim-1}
Let $X$ be a proper scheme of dimension $\leq 1$ over a field $k$.
There exists a dualizing complex $\omega_X^\bullet$ with the
following properties
\begin{enumerate}
\item $H^i(\omega_X^\bullet)$ is nonzero only for $i = -1, 0$,
\item $\omega_X = H^{-1}(\omega_X^\bullet)$
is a coherent Cohen-Macaulay module whose support is the
irreducible components of dimension $1$,
\item for $x \in X$ closed, the module $H^0(\omega_{X, x}^\bullet)$
is nonzero if and only if either
\begin{enumerate}
\item $\dim(\mathcal{O}_{X, x}) = 0$ or
\item $\dim(\mathcal{O}_{X, x}) = 1$
and $\mathcal{O}_{X, x}$ is not Cohen-Macaulay,
\end{enumerate}
\item for $K \in D_\QCoh(\mathcal{O}_X)$ there are functorial
isomorphisms\footnote{This property
characterizes $\omega_X^\bullet$ in $D_\QCoh(\mathcal{O}_X)$
up to unique isomorphism by the Yoneda lemma. Since $\omega_X^\bullet$
is in $D^b_{\textit{Coh}}(\mathcal{O}_X)$ in fact it suffices to consider
$K \in D^b_{\textit{Coh}}(\mathcal{O}_X)$.}
$$
\Ext^i_X(K, \omega_X^\bullet) = \Hom_k(H^{-i}(X, K), k)
$$
compatible with shifts and distinguished triangles,
\item there are functorial isomorphisms
$\Hom(\mathcal{F}, \omega_X) = \Hom_k(H^1(X, \mathcal{F}), k)$
for $\mathcal{F}$ quasi-coherent on $X$,
\item if $X \to \Spec(k)$ is smooth of relative dimension $1$,
then $\omega_X \cong \Omega_{X/k}$.
\end{enumerate}
\end{lemma}

\begin{lemma}
\label{lemma-global-generation}
Let $k$ be a field. Let $X$ be a proper scheme over $k$
which is reduced, connected, and of dimension $1$.
Let $\mathcal{L}$ be an invertible $\mathcal{O}_X$-module.
Assume that for every reduced connected closed subscheme
$Y \subset X$ of dimension $1$ we have
$$
\deg(\mathcal{L}|_Y) \geq 2\dim_k H^1(Y, \mathcal{O}_Y)
$$
Then $\mathcal{L}$ is globally generated.
\end{lemma}

\begin{proof}
By induction on the number of irreducible components of $X$.
If $X$ is irreducible, then the lemma holds by
Lemma \href{https://stacks.math.columbia.edu/tag/0E3C}{lemma degree more than 2g (Tag 0E3C)}
applied to $X$ viewed as a scheme over the field
$k' = H^0(X, \mathcal{O}_X)$. Assume $X$ is not irreducible.
Before we continue, if $k$ is finite, then we replace $k$
by a purely transcendental extension $K$. This is allowed by
Varieties, Lemmas
\href{https://stacks.math.columbia.edu/tag/0B57}{lemma globally generated base change (Tag 0B57)},
\href{https://stacks.math.columbia.edu/tag/0B59}{lemma degree base change (Tag 0B59)},
\href{https://stacks.math.columbia.edu/tag/035Z}{lemma geometrically reduced any base change (Tag 035Z)}, and
\href{https://stacks.math.columbia.edu/tag/038F}{lemma bijection irreducible components (Tag 038F)},
Cohomology of Schemes, Lemma \href{https://stacks.math.columbia.edu/tag/02KH}{lemma flat base change cohomology (Tag 02KH)},
Lemma \href{https://stacks.math.columbia.edu/tag/0E32}{lemma sanity check duality (Tag 0E32)} and the elementary fact that
$K$ is geometrically integral over $k$.

\medskip\noindent
Assume that $\mathcal{L}$ is not globally generated to get a contradiction.
Then we may choose a coherent ideal sheaf
$\mathcal{I} \subset \mathcal{O}_X$ such that
$H^0(X, \mathcal{I}\mathcal{L}) = H^0(X, \mathcal{L})$
and such that $\mathcal{O}_X/\mathcal{I}$ is nonzero with
support of dimension $0$. For example, take $\mathcal{I}$
the ideal sheaf of any closed point in the common
vanishing locus of the global sections of $\mathcal{L}$.
We consider the short exact sequence
$$
0 \to \mathcal{I}\mathcal{L} \to \mathcal{L} \to 
\mathcal{L}/\mathcal{I}\mathcal{L} \to 0
$$
Since the support of $\mathcal{L}/\mathcal{I}\mathcal{L}$
has dimension $0$ we see that $\mathcal{L}/\mathcal{I}\mathcal{L}$
is generated by global sections
(Varieties, Lemma \href{https://stacks.math.columbia.edu/tag/0AYT}{lemma chi tensor finite (Tag 0AYT)}).
From the short exact sequence,
and the fact that $H^0(X, \mathcal{I}\mathcal{L}) = H^0(X, \mathcal{L})$
we get an injection
$H^0(X, \mathcal{L}/\mathcal{I}\mathcal{L}) \to H^1(X, \mathcal{I}\mathcal{L})$.

\medskip\noindent
Recall that the $k$-vector space $H^1(X, \mathcal{I}\mathcal{L})$
is dual to $\Hom(\mathcal{I}\mathcal{L}, \omega_X)$.
Choose $\varphi : \mathcal{I}\mathcal{L} \to \omega_X$.
By Lemma \href{https://stacks.math.columbia.edu/tag/0E3E}{lemma vanishing on gorenstein (Tag 0E3E)} we have
$H^1(X, \mathcal{L}) = 0$. Hence
$$
\dim_k H^0(X, \mathcal{I}\mathcal{L}) = \dim_k H^0(X, \mathcal{L}) =
\deg(\mathcal{L}) + \chi(\mathcal{O}_X) > \dim_k H^1(X, \mathcal{O}_X) =
\dim_k H^0(X, \omega_X)
$$
We conclude that $\varphi$ is not injective on global sections, in particular
$\varphi$ is not injective. For every generic point $\eta \in X$
of an irreducible component of $X$ denote
$V_\eta \subset \Hom(\mathcal{I}\mathcal{L}, \omega_X)$ the $k$-subvector
space consisting of those $\varphi$ which are zero at $\eta$.
Since every associated point of $\mathcal{I}\mathcal{L}$
is a generic point of $X$, the above shows that
$\Hom(\mathcal{I}\mathcal{L}, \omega_X) = \bigcup V_\eta$.
As $X$ has finitely many generic points and $k$ is infinite, we conclude
$\Hom(\mathcal{I}\mathcal{L}, \omega_X) = V_\eta$ for some $\eta$.
Let $\eta \in C \subset X$ be the corresponding irreducible component.
Let $Y \subset X$ be the union of the other irreducible components
of $X$. Then $Y$ is a nonempty reduced closed subscheme not equal to $X$.
Let $\mathcal{J} \subset \mathcal{O}_X$ be the ideal sheaf of $Y$.
Please keep in mind that the support of $\mathcal{J}$ is $C$.

\medskip\noindent
Let $\varphi : \mathcal{I}\mathcal{L} \to \omega_X$ be arbitrary.
Since $\mathcal{J}\mathcal{I}\mathcal{L}$
has no embedded associated points
(as a submodule of $\mathcal{L}$) and as $\varphi$ is zero
in the generic point $\eta$ of the support of $\mathcal{J}$, we find that
$\varphi$ factors as
$$
\mathcal{I}\mathcal{L} \to
\mathcal{I}\mathcal{L}/\mathcal{J}\mathcal{I}\mathcal{L} \to \omega_X
$$
We can view $\mathcal{I}\mathcal{L}/\mathcal{J}\mathcal{I}\mathcal{L}$
as the pushforward of a coherent sheaf on $Y$ which by abuse of
notation we indicate with the same symbol.
Since $\omega_Y = \SheafHom(\mathcal{O}_Y, \omega_X)$
by Lemma \href{https://stacks.math.columbia.edu/tag/0E33}{lemma closed immersion dim 1 CM (Tag 0E33)}
we find a factorization
$$
\mathcal{I}\mathcal{L} \to
\mathcal{I}\mathcal{L}/
\mathcal{J}\mathcal{I}\mathcal{L} 
\xrightarrow{\varphi'} \omega_Y \to \omega_X
$$
of $\varphi$. Let $\mathcal{I}' \subset \mathcal{O}_Y$ be the
image of $\mathcal{I} \subset \mathcal{O}_X$. There is a surjective map
$\mathcal{I}\mathcal{L}/\mathcal{J}\mathcal{I}\mathcal{L} \to
\mathcal{I}'\mathcal{L}|_Y$ whose kernel is supported in closed points.
Since $\omega_Y$ is a Cohen-Macaulay module on $Y$, the map $\varphi'$
factors through a map
$\varphi'' : \mathcal{I}'\mathcal{L}|_Y \to \omega_Y$.
Thus we have commutative diagrams
$$
\vcenter{
\xymatrix{
0 \ar[r] &
\mathcal{I}\mathcal{L} \ar[r] \ar[d] &
\mathcal{L} \ar[r] \ar[d] &
\mathcal{L}/\mathcal{I}\mathcal{L} \ar[r] \ar[d] &
0 \\
0 \ar[r] &
\mathcal{I}'\mathcal{L}|_Y \ar[r] &
\mathcal{L}|_Y \ar[r] &
\mathcal{L}|_Y/\mathcal{I}'\mathcal{L}|_Y \ar[r] &
0
}
}
\quad\text{and}\quad
\vcenter{
\xymatrix{
\mathcal{I}\mathcal{L} \ar[r]_\varphi \ar[d] & \omega_X \\
\mathcal{I}'\mathcal{L}|_Y \ar[r]^{\varphi''} & \omega_Y \ar[u]
}
}
$$
Now we can finish the proof as follows:
Since for every $\varphi$ we have a $\varphi''$ and since
$\omega_X \in \textit{Coh}(\mathcal{O}_X)$
represents the functor $\mathcal{F} \mapsto \Hom_k(H^1(X, \mathcal{F}), k)$,
we find that
$H^1(X, \mathcal{I}\mathcal{L}) \to H^1(Y, \mathcal{I}'\mathcal{L}|_Y)$
is injective. Since the boundary
$H^0(X, \mathcal{L}/\mathcal{I}\mathcal{L}) \to H^1(X, \mathcal{I}\mathcal{L})$
is injective, we conclude the composition
$$
H^0(X, \mathcal{L}/\mathcal{I}\mathcal{L}) \to
H^0(X, \mathcal{L}|_Y/\mathcal{I}'\mathcal{L}|_Y) \to
H^1(X, \mathcal{I}'\mathcal{L}|_Y)
$$
is injective. Since
$\mathcal{L}/\mathcal{I}\mathcal{L} \to
\mathcal{L}|_Y/\mathcal{I}'\mathcal{L}|_Y$
is a surjective map of coherent modules whose supports have
dimension $0$, we see that the first map
$H^0(X, \mathcal{L}/\mathcal{I}\mathcal{L}) \to
H^0(X, \mathcal{L}|_Y/\mathcal{I}'\mathcal{L}|_Y)$
is surjective (and hence bijective).
But by induction we have that $\mathcal{L}|_Y$ is globally
generated (if $Y$ is disconnected this still works of course)
and hence the boundary map
$$
H^0(X, \mathcal{L}|_Y/\mathcal{I}'\mathcal{L}|_Y) \to
H^1(X, \mathcal{I}'\mathcal{L}|_Y)
$$
cannot be injective.
This contradiction finishes the proof.
\end{proof}
\end{document}
